\chapter{Integrated Correlators and Bootstrap on Conformal Manifolds}
\label{ch:volume-iii-integrated-correlators-bootstrap-conformal-manifolds}

\section{Integrated Correlators as Linear Functionals}

Let \(p\in\mathcal M_{\rm conf}\) be a point on a conformal manifold and let
\(\mathcal O_a\) be an exactly marginal operator.  Differentiating a
renormalized generating functional with respect to \(\lambda^a\) inserts
\(\int\mathcal O_a\).  For a set of local sources \(J_i(x)\), write
\[
  W[J,\lambda]=\log Z[J,\lambda].
\]
Then
\begin{equation}
  \partial_a
  \left.
  \frac{\delta^n W}{\delta J_1(x_1)\cdots\delta J_n(x_n)}
  \right|_{J=0}
  =
  -\int\dd^Dx\,
  \left\langle
    \mathcal O_a(x)\mathcal X_1(x_1)\cdots\mathcal X_n(x_n)
  \right\rangle_{\rm conn}
  +\text{contact terms}.
  \label{eq:coupling-derivative-integrated-insertion}
\end{equation}
The contact terms express the renormalization prescription for collisions of
\(\mathcal O_a(x)\) with the other insertions and with itself under repeated
derivatives.

On compact supersymmetric backgrounds, localization often computes quantities
of the form
\[
  \mathcal I[\mathcal G]
  =
  \int \dd\mu(x_i)\,K(x_i)\,\mathcal G(x_i),
\]
where \(\mathcal G\) is a flat-space CFT correlator in a protected multiplet,
\(K(x_i)\) is a known kernel fixed by the background and the supersymmetric
deformation, and \(\dd\mu\) is the conformally invariant integration measure
with the appropriate subtractions.  Thus localization can produce exact
linear functionals of CFT correlators.

\begin{figure}[H]
  \centering
  \resizebox{0.92\textwidth}{!}{%
  \begin{tikzpicture}[x=1cm,y=1cm,every node/.style={font=\scriptsize}]
    \tikzset{
      arrow/.style={-{Latex[length=2mm]},line width=0.85pt},
      box/.style={draw,rounded corners=4pt,line width=0.8pt,fill=blue!4,
        minimum width=3.1cm,minimum height=0.9cm,align=center}
    }
    \node[box] (loc) at (-4.2,0.8) {localized\\observable};
    \node[box] (der) at (0,0.8) {coupling or mass\\derivatives};
    \node[box] (int) at (4.2,0.8) {integrated\\correlator};
    \node[box] (kernel) at (0,-0.8) {known kernel\\\(K(x_i)\)};
    \node[box] (data) at (4.2,-0.8) {constraint on\\CFT data};
    \draw[arrow] (loc)--(der);
    \draw[arrow] (der)--(int);
    \draw[arrow] (kernel)--(int);
    \draw[arrow] (int)--(data);
    \node[green!45!black,align=center] at (-4.0,-0.85)
      {exact protected computation\\becomes a bootstrap input};
  \end{tikzpicture}
  }
  \caption{Derivatives of exact observables yield integrated correlators,
  which are linear constraints on the ordinary CFT data.}
  \label{fig:integrated-correlator-functional}
\end{figure}

\section{Differentiating Crossing}

For identical scalar primaries \(\phi\), the crossing equation can be written
as
\begin{equation}
  \sum_{\mathcal X}
  \lambda_{\phi\phi\mathcal X}(p)^2
  F_{\Delta_{\mathcal X}(p),\ell_{\mathcal X}}(u,v)=0,
  \label{eq:crossing-family-pointwise}
\end{equation}
where \(F_{\Delta,\ell}\) is the crossing vector built from conformal blocks
and the identity term is included in the sum with its fixed coefficient.
Along a conformal manifold, the CFT data depend on \(p\).  Differentiating
\eqref{eq:crossing-family-pointwise} in a tangent direction \(a\) gives the
linearized crossing equation
\begin{equation}
  \sum_{\mathcal X}
  \left[
    2\lambda_{\phi\phi\mathcal X}
      \nabla_a\lambda_{\phi\phi\mathcal X}\,
      F_{\Delta_{\mathcal X},\ell_{\mathcal X}}
    +
    \lambda_{\phi\phi\mathcal X}^2
      (\nabla_a\Delta_{\mathcal X})
      \partial_\Delta F_{\Delta_{\mathcal X},\ell_{\mathcal X}}
  \right]
  =
  0.
  \label{eq:linearized-crossing-conformal-manifold}
\end{equation}
Here \(\nabla_a\) denotes the connection on the relevant operator bundles.  In
degenerate sectors, the scalar quantities in
\eqref{eq:linearized-crossing-conformal-manifold} become matrices, and the
equation is written after choosing a local orthonormal frame.

Higher derivatives of crossing produce higher-order constraints involving
\(\nabla_a\nabla_b\Delta_{\mathcal X}\), derivatives of OPE coefficients, and
contact terms from repeated integrated insertions.  The resulting hierarchy is
the conformal-perturbation expansion of crossing symmetry.

\begin{figure}[H]
  \centering
  \resizebox{0.92\textwidth}{!}{%
  \begin{tikzpicture}[x=1cm,y=1cm,every node/.style={font=\scriptsize}]
    \tikzset{
      arrow/.style={-{Latex[length=2mm]},line width=0.85pt},
      curve/.style={blue!65!black,line width=1.0pt},
      box/.style={draw,rounded corners=4pt,fill=blue!4,line width=0.8pt,
        minimum width=3.2cm,minimum height=0.82cm,align=center}
    }
    \draw[curve] (-4.2,-0.45)
      .. controls (-2.7,0.75) and (-1.2,0.75) .. (0.2,-0.15)
      .. controls (1.6,-1.0) and (3.0,-0.55) .. (4.2,0.38);
    \node[blue!65!black] at (2.55,1.0) {\(\mathcal M_{\rm conf}\)};
    \foreach \x/\y/\lab in {-2.4/0.42/{p},-0.55/0.03/{p+\delta p},1.55/-0.48/{p+2\delta p}} {
      \fill[violet!80!black] (\x,\y) circle (2pt);
      \node[violet!80!black,above] at (\x,\y) {\(\lab\)};
    }
    \node[box] at (-2.55,-1.25) {crossing\\at \(p\)};
    \node[box] at (0.0,-1.85) {linearized\\crossing};
    \node[box] at (2.85,-1.25) {higher\\derivatives};
    \draw[arrow] (-1.62,-1.35)--(-1.05,-1.62);
    \draw[arrow] (1.15,-1.62)--(1.90,-1.35);
  \end{tikzpicture}
  }
  \caption{Crossing holds at every point of a conformal manifold.  Its
  derivatives constrain the variation of dimensions and OPE coefficients.}
  \label{fig:crossing-differentiated-on-manifold}
\end{figure}

\section{Protected Anchors}

Let \(\mathcal P(p)\) denote a set of protected data at \(p\): protected
dimensions, protected OPE coefficients, an index, a chiral-algebra vacuum
character, a sphere partition function, or a Wilson-loop expectation value.
Each element of \(\mathcal P(p)\) is a function of the full CFT data:
\[
  \mathcal P_A(p)
  =
  \mathcal F_A\!\left[
    \Delta_i(p),C_{ijk}(p),\ldots
  \right].
\]
When \(\mathcal P_A(p)\) is known exactly, it becomes a constraint on the
ordinary crossing problem.

For example, suppose an exact integrated correlator is known:
\[
  \mathcal I[\mathcal G]=\mathcal C_{\rm exact}(\lambda,\bar\lambda).
\]
Expanding \(\mathcal G\) in conformal blocks gives
\begin{equation}
  \sum_{\mathcal X}
  \lambda_{\phi\phi\mathcal X}^2
  \mathcal I\!\left[
    G_{\Delta_{\mathcal X},\ell_{\mathcal X}}
  \right]
  =
  \mathcal C_{\rm exact}(\lambda,\bar\lambda).
  \label{eq:integrated-correlator-block-constraint}
\end{equation}
The numbers
\[
  \mathcal I\!\left[
    G_{\Delta,\ell}
  \right]
\]
are computable kinematic functions.  Equation
\eqref{eq:integrated-correlator-block-constraint} is therefore a linear
constraint on positive OPE data whenever the relevant OPE coefficients enter
as squares.

\section{Duality and Automorphic Constraints}

Let a duality group \(\Gamma\) act on the coordinate cover
\(\widetilde{\mathcal M}_{\rm conf}\).  A protected observable
\(\mathcal P(p)\) may transform with a weight or with a finite-dimensional
matrix:
\[
  \mathcal P(\gamma p)
  =
  \rho(\gamma,p)\,\mathcal P(p).
\]
Here \(\rho(\gamma,p)\) is determined by the normalization of the observable
and by the duality action on background fields and operator bundles.

The same covariance applies to unprotected data after operator labels are
transported:
\[
  \Delta_i(\gamma p)=\Delta_{\gamma\cdot i}(p),
  \qquad
  C_{ijk}(\gamma p)
  =
  M_i{}^{i'}M_j{}^{j'}M_k{}^{k'}C_{i'j'k'}(p).
\]
Combining crossing with duality produces a global bootstrap problem on
\(\mathcal M_{\rm conf}\).  Local perturbative data near a weak-coupling cusp,
strong-coupling expansions from a dual frame, protected exact functions, and
positivity constraints must fit into a single set of globally well-defined
CFT data.

\begin{figure}[H]
  \centering
  \resizebox{0.92\textwidth}{!}{%
  \begin{tikzpicture}[x=1cm,y=1cm,every node/.style={font=\scriptsize}]
    \tikzset{
      arrow/.style={-{Latex[length=2mm]},line width=0.85pt},
      box/.style={draw,rounded corners=4pt,fill=blue!4,line width=0.8pt,
        minimum width=3.0cm,minimum height=0.85cm,align=center}
    }
    \node[box] (cross) at (-4.2,1.35) {crossing and\\positivity};
    \node[box] (dual) at (-4.2,0.45) {duality\\covariance};
    \node[box] (prot) at (-4.2,-0.45) {protected exact\\observables};
    \node[box] (pert) at (-4.2,-1.35) {local perturbative\\expansions};
    \node[box] (global) at (3.1,0) {global CFT data\\over
      \(\mathcal M_{\rm conf}\)};
    \draw[arrow] (cross.east)--(global.west);
    \draw[arrow] (dual.east)--(global.west);
    \draw[arrow] (prot.east)--(global.west);
    \draw[arrow] (pert.east)--(global.west);
    \node[green!45!black,align=center] at (-0.45,-2.12)
      {the bootstrap problem becomes a global compatibility problem};
  \end{tikzpicture}
  }
  \caption{On a conformal manifold, exact protected quantities, duality,
  perturbation theory, and crossing constrain one global set of CFT data.}
  \label{fig:global-bootstrap-conformal-manifold}
\end{figure}

\section{Large N Integrated Constraints}

For a holographic large \(N\) CFT, write a connected four-point function as
\[
  \mathcal G
  =
  \mathcal G_{\rm disc}
  +
  C_T^{-1}\mathcal G^{(1)}
  +
  C_T^{-2}\mathcal G^{(2)}
  +\cdots .
\]
At order \(C_T^{-1}\), \(\mathcal G^{(1)}\) is computed by tree-level Witten
diagrams.  In Mellin space, a local ansatz is a sum of exchange terms and
crossing-symmetric polynomials:
\[
  M^{(1)}(s,t)
  =
  M_{\rm exch}(s,t)
  +
  \sum_{m,n} a_{mn}\,s^m t^n
\]
with the sum restricted by the derivative order.  An exact integrated
correlator gives
\[
  \mathcal I[M^{(1)}]=\mathcal C^{(1)}_{\rm exact},
\]
which fixes a linear combination of the coefficients \(a_{mn}\).  A family of
integrated constraints can determine higher-derivative bulk couplings that
would otherwise be invisible to symmetry alone.

The same logic applies at higher orders, with loop Witten diagrams and
multi-trace anomalous dimensions replacing the tree-level data.  Integrated
correlators are therefore a precise interface between localization, conformal
bootstrap, and holographic effective field theory.
